\section{Security}
\label{sec:security}
The system operates on a university research lab LAN and is not exposed to the public internet; accordingly, the security measures are pragmatic rather than exhaustive. Security was planned in three phases. As built, the system is at an earlier stage than Phase~1.

\paragraph{Current State}
\begin{itemize}[noitemsep]
  \item The broker accepts anonymous connections; nodes connect without credentials.
  \item Exposure is limited by the Docker port mapping: by default the broker is reachable only from the historian server itself, and is opened to the lab LAN only when real nodes are connected.
\end{itemize}

\paragraph{Phase 1: Authentication (deferred)}
\begin{itemize}[noitemsep]
  \item Disable anonymous connections.
  \item Enable username and password authentication via Mosquitto's built-in password file mechanism.
  \item Issue separate credentials per client (one per DAQ node, one for the Telegraf bridge).
\end{itemize}

\paragraph{Phase 2: Hardening (deferred)} ~\\
Access control lists (ACLs) enforce the principle of least privilege. Each DAQ node may only publish to its own topic subtree, so that a compromised node cannot spoof measurements from another node. The Telegraf bridge may subscribe to the full topic tree but cannot publish.

\paragraph{Phase 3: TLS (future)} ~\\
If the broker is ever exposed outside the lab VLAN, mutual TLS on port 8883 with certificates issued by a local certificate authority is the intended upgrade path. This is not required for the current testbed scope.
